%% =======================================================
%% Actividad X:
%% 
%% 
%% =======================================================

%% ---- CONFIGURACIÓN ------------------------

\documentclass[11pt,letterpaper]{article}

\usepackage[sinindice]{tecnm}
\usepackage[T1]{fontenc}
\usepackage{xcolor}
\usepackage[most]{tcolorbox} % Carga tcolorbox con soporte para listings y saltos de página
\tcbuselibrary{listings, skins, breakable}

% Paleta de colores (estilo GitHub Light / VS Code)
\definecolor{codebg}{HTML}{F8F9FA}       % Fondo del código
\definecolor{codehead}{HTML}{E9ECEF}     % Fondo de la barra de título
\definecolor{codeborder}{HTML}{CED4DA}   % Borde de la tarjeta
\definecolor{codetitle}{HTML}{212529}    % Texto del título
\definecolor{codekw}{HTML}{0550AE}       % Palabras clave (import, def, for...)
\definecolor{codefunc}{HTML}{8250DF}     % Funciones y métodos (describe, groupby...)
\definecolor{codestr}{HTML}{0A3069}      % Strings
\definecolor{codecomm}{HTML}{6E7781}     % Comentarios
\definecolor{codenum}{HTML}{8C959F}      % Números de línea
\definecolor{codeid}{HTML}{24292F}       % Texto general

% Configuración base de listings
\lstdefinestyle{pythonestilo}{
  language=Python,
  basicstyle=\ttfamily\footnotesize\color{codeid},
  keywordstyle=\color{codekw}\bfseries,
  keywordstyle=[2]\color{codefunc}, % Estilo secundario para métodos de datos
  stringstyle=\color{teal!80!black},
  commentstyle=\color{codecomm}\itshape,
  morekeywords={as, True, False, None, with, yield, lambda},
  morekeywords=[2]{read_csv, head, tail, info, describe, groupby, mean, value_counts, to_csv, print, len, range},
  numbers=left,
  numberstyle=\tiny\ttfamily\color{codenum},
  stepnumber=1,
  numbersep=10pt,
  showstringspaces=false,
  morestring=[b]{"""},
  tabsize=4,
  breaklines=true,
  breakatwhitespace=true,
  columns=flexible,
  keepspaces=true,
  literate=
    {á}{{\'a}}1 {é}{{\'e}}1 {í}{{\'i}}1 {ó}{{\'o}}1 {ú}{{\'u}}1
    {Á}{{\'A}}1 {É}{{\'E}}1 {Í}{{\'I}}1 {Ó}{{\'O}}1 {Ú}{{\'U}}1
    {ñ}{{\~n}}1 {Ñ}{{\~N}}1 {¿}{{?`}}1 {¡}{{!`}}1
}

% Definición del entorno personalizado \begin{codigo}{Título}{etiqueta}
% \newtcblisting[auto counter, number within=section]{codigo}[2][]{
\newtcblisting[auto counter]{codigo}[2][]{
  listing only,
  breakable,                    % Permite que el bloque se divida si pasa de página
  enhanced,
  colback=codebg,
  colframe=codeborder,
  colbacktitle=codehead,
  coltitle=codetitle,
  fonttitle=\bfseries\small\sffamily,
  boxrule=0.6pt,
  arc=4pt,                      % Bordes redondeados
  left=20pt,                    % Espacio izquierdo para los números de línea
  right=8pt,
  top=6pt,
  bottom=6pt,
  toptitle=4pt,
  bottomtitle=4pt,
  before skip=12pt,
  after skip=12pt,
  listing options={style=pythonestilo},
  title={Código \thetcbcounter: #2},
  #1
}

% Paleta para salida de consola (estilo terminal limpia)
\definecolor{termbg}{HTML}{212529}      % Fondo oscuro tipo terminal
\definecolor{termhead}{HTML}{343A40}    % Barra superior
\definecolor{termtext}{HTML}{F8F9FA}    % Texto claro de alta legibilidad
\definecolor{termborder}{HTML}{495057}  % Borde

% Entorno \begin{consola}[label=out:iris]{Título}
\newtcblisting[auto counter]{consola}[2][]{
  listing only,
  breakable,
  colback=termbg,
  colframe=termborder,
  colbacktitle=termhead,
  coltitle=termtext,
  colupper=termtext,
  fonttitle=\bfseries\small\sffamily,
  boxrule=0.6pt,
  arc=4pt,
  left=8pt,
  right=8pt,
  top=6pt,
  bottom=6pt,
  toptitle=4pt,
  bottomtitle=4pt,
  before skip=12pt,
  after skip=12pt,
  listing options={
    basicstyle=\ttfamily\scriptsize\color{termtext}, % \scriptsize evita que las tablas de Pandas se partan de línea
    breaklines=true,
    columns=fixed,        % Mantiene alineadas las columnas de Pandas
    keepspaces=true,
    showstringspaces=false,
    literate=
      {á}{{\'a}}1 {é}{{\'e}}1 {í}{{\'i}}1 {ó}{{\'o}}1 {ú}{{\'u}}1
      {Á}{{\'A}}1 {É}{{\'E}}1 {Í}{{\'I}}1 {Ó}{{\'O}}1 {Ú}{{\'U}}1
      {ñ}{{\~n}}1 {Ñ}{{\~N}}1
  },
  title={Salida \thetcbcounter: #2},
  #1
}

%% ---- DATOS DE LA MATERIA ------------------------

\materia{Modelos Estadísticos para el Aprendizaje \\ Automático y Ciencia de Datos}
\actividad{Tema 1.1
Introducción a Conjuntos de Datos Reales y Abiertos \\ Actividad 1 - Proyecto Explorador Avanzado del Dataset Iris}

\profesor{Dr. José Ismael Ojeda Campaña}
%\fecha{15 de octubre de 2026}    % opcional; por defecto es la fecha de hoy

\begin{document}
\portada
\indicegeneral

%% ---- CONTENIDO ---------------------------------

\section{Introducción}

El conjunto de datos \textit{Iris} \parencite{fisher1936}  es uno de los referentes más icónicos en el aprendizaje automático \parencite{meaa_u1_tema11}  y la estadística multivariada, constituyendo un punto de partida estándar para validar flujos de trabajo en ciencia de datos. A través de sus 150 observaciones y cuatro variables morfológicas medidas en centímetros (longitud y ancho de sépalos y pétalos), permite evaluar técnicas de análisis exploratorio, agrupación y clasificación supervisada entre tres especies: \textit{Iris setosa}, \textit{Iris versicolor} e \textit{Iris virginica}.

\section{Desarrollo de la Actividad}

\textbf{Objetivo:} Extender el código base proporcionado para ejecutar un Análisis Exploratorio de Datos (EDA) completo sobre el dataset Iris, incorporando segmentación por especies y exportación automatizada de resultados.

\subsection*{Requerimientos del Proyecto}

\begin{description}
    \item[Tarea 1: Extensión del Código Base.] Modificar el script inicial para inspeccionar la estructura y distribución general de los datos:
    \begin{enumerate}
        \item Visualizar las primeras 10 filas del conjunto de datos mediante \texttt{df.head(10)}.
        \item Visualizar las últimas 8 filas mediante \texttt{df.tail(8)}.
        \item Obtener la estructura general, tipos de datos y conteo de valores no nulos con \texttt{df.info()}.
        \item Calcular las estadísticas descriptivas globales (media, desviación estándar, cuartiles) con \texttt{df.describe()}.
    \end{enumerate}

    \item[Tarea 2: Análisis por Especies.] Segmentar la información según la variable objetivo para identificar patrones morfológicos:
    \begin{enumerate}
        \item Mostrar la distribución y número de muestras por cada especie.
        \item Calcular los valores promedio de cada característica agrupados por especie.
        \item Generar un resumen estadístico completo agrupado por especie.
    \end{enumerate}

    \item[Tarea 3: Exportación de Resultados.] Automatizar la persistencia de los datos procesados y los hallazgos obtenidos:
    \begin{enumerate}
        \item Exportar el \textit{DataFrame} completo a un archivo en formato CSV.
        \item Almacenar las tablas estadísticas en un archivo de texto plano (\texttt{.txt}).
        \item Generar de forma automática un reporte resumen con los principales hallazgos del análisis.
    \end{enumerate}
\end{description}

\section {Desarrollo}

Se desarrollaron las actividades anteriores en el Código \ref{cod:iris} que puede verse a continuación:

\begin{codigo}[label=cod:iris]{Exploración Avanzada del Dataset Iris}

import pandas as pd
from sklearn.datasets import load_iris

# Actividad 1: Explorador Avanzado del Dataset Iris
# Alumno: Fernando Valderrábano Reyes
# Asignatura: Modelos estadísticos para el aprendizaje automático y ciencia de datos
# Docente: Dr. José Ismael Ojeda Campaña

"""
Tarea 1: Extensión del Código Base
Modificar el código para mostrar:
a) Las primeras 10 filas
b) Las últimas 8 filas
c) Información general del dataset con df.info()
d) Estadísticas descriptivas con df.describe()
"""

# Cargar el conjunto de datos Iris
data = load_iris()
df = pd.DataFrame(data.data, columns=data.feature_names)
df['species'] = pd.Categorical.from_codes(data.target, data.target_names)

print("\nPrimeras 10 filas del conjunto de datos:")
print(df.head(10))
print ("\n");

print("\nÚltimas 8 filas del conjunto de datos:")
print(df.tail(8))
print ("\n");

print("\nInformación general del conjunto de datos:")
print(df.info())
print ("\n");

print("\nEstadisticas descriptivas del conjunto de datos:")
print(df.describe())
print ("\n");

"""
Tarea 2: Análisis por Especies
Agregar código para:
a) Mostrar el número de muestras por especie
b) Calcular promedios de cada característica por especie
c) Generar un resumen estadístico agrupado por especie
"""

"""Mostrar el número de muestras por especie"""
print("\nNúmero de muestras por especie:")
print(df['species'].value_counts())
print ("\n");

"""Calcular promedios de cada característica por especie"""
print("\nPromedios de cada característica por especie:")
print(df.groupby('species').mean())
print ("\n");

"""Generar un resumen estadístico agrupado por especie"""
print("\nResumen estadístico agrupado por especie:")
print(df.groupby('species').describe())

"""
Tarea 3: Exportación de Resultados
Implementar:
a) Exportar el DataFrame completo a CSV
b) Guardar estadísticas en un archivo de texto
c) Crear un reporte automático con los hallazgos
"""

# Exportar el DataFrame completo a CSV
df.to_csv('iris_dataset.csv', index=False)

# Guardar estadísticas en un archivo de texto
with open('iris_statistics.txt', 'w') as f:
    f.write("Número de muestras por especie:\n")
    f.write(str(df['species'].value_counts()))
    f.write("\n\nPromedios de cada característica por especie:\n")
    f.write(str(df.groupby('species').mean()))
    f.write("\n\nResumen estadístico agrupado por especie:\n")
    f.write(str(df.groupby('species').describe()))

# Crear un reporte automático con los hallazgos
with open('iris_report.txt', 'w') as f:
    f.write("Reporte del Conjunto de Datos Iris\n")
    f.write("="*40 + "\n\n")
    f.write("Primeras 10 filas del conjunto de datos:\n")
    f.write(str(df.head(10)))
    f.write("\n\nÚltimas 8 filas del conjunto de datos:\n")
    f.write(str(df.tail(8)))
    f.write("\n\nInformación general del conjunto de datos:\n")
    f.write(str(df.info()))
    f.write("\n\nEstadísticas descriptivas del conjunto de datos:\n")
    f.write(str(df.describe()))
    f.write("\n\nNúmero de muestras por especie:\n")
    f.write(str(df['species'].value_counts()))
    f.write("\n\nPromedios de cada característica por especie:\n")
    f.write(str(df.groupby('species').mean()))
    f.write("\n\nResumen estadístico agrupado por especie:\n")
    f.write(str(df.groupby('species').describe()))

\end{codigo}


\subsection*{Resultados}

Se implementaron satisfactoriamente las actividades correspondientes a las Tareas 1, 2 y 3, cuyo código fuente íntegro se presenta en el Código~\ref{cod:iris}. El cuaderno interactivo puede consultarse en Google Colab:

\begin{center}
    \url{https://colab.research.google.com/drive/1WpdbhrKdbbJqNeG-X2yIbUvCkw1x1dhJ?usp=sharing}
\end{center}

En la Tabla~\ref{tab:promedios_iris} se resumen los valores promedio obtenidos por especie (Tarea 2), donde se observa una clara separación morfológica en las dimensiones del pétalo entre \textit{Iris setosa} y las demás especies.

\begin{table}[htbp]
\centering
\caption{Promedios de las características morfológicas (en cm) agrupados por especie.}
\label{tab:promedios_iris}
\begin{tabular}{lcccc}
\hline
\textbf{Especie} & \textbf{Longitud Sépalo} & \textbf{Ancho Sépalo} & \textbf{Longitud Pétalo} & \textbf{Ancho Pétalo} \\
\hline
\textit{Iris setosa}     & 5.006 & 3.428 & 1.462 & 0.246 \\
\textit{Iris versicolor} & 5.936 & 2.770 & 4.260 & 1.326 \\
\textit{Iris virginica}  & 6.588 & 2.974 & 5.552 & 2.026 \\
\hline
\end{tabular}
\end{table}

Asimismo, la salida completa generada en consola durante la ejecución del script (inspección de registros, tipos de datos y estadísticas descriptivas) se muestra en la Salida~\ref{out:iris}.

\begin{consola}[label=out:iris]{Ejecución en terminal de \texttt{actividad001.py}}
Primeras 10 filas del conjunto de datos:
   sepal length (cm)  sepal width (cm)  petal length (cm)  petal width (cm) species
0                5.1               3.5                1.4               0.2  setosa
1                4.9               3.0                1.4               0.2  setosa
2                4.7               3.2                1.3               0.2  setosa
3                4.6               3.1                1.5               0.2  setosa
4                5.0               3.6                1.4               0.2  setosa
5                5.4               3.9                1.7               0.4  setosa
6                4.6               3.4                1.4               0.3  setosa
7                5.0               3.4                1.5               0.2  setosa
8                4.4               2.9                1.4               0.2  setosa
9                4.9               3.1                1.5               0.1  setosa



Últimas 8 filas del conjunto de datos:
     sepal length (cm)  sepal width (cm)  petal length (cm)  petal width (cm)    species
142                5.8               2.7                5.1               1.9  virginica
143                6.8               3.2                5.9               2.3  virginica
144                6.7               3.3                5.7               2.5  virginica
145                6.7               3.0                5.2               2.3  virginica
146                6.3               2.5                5.0               1.9  virginica
147                6.5               3.0                5.2               2.0  virginica
148                6.2               3.4                5.4               2.3  virginica
149                5.9               3.0                5.1               1.8  virginica



Información general del conjunto de datos:
<class 'pandas.DataFrame'>
RangeIndex: 150 entries, 0 to 149
Data columns (total 5 columns):
 #   Column             Non-Null Count  Dtype   
---  ------             --------------  -----   
 0   sepal length (cm)  150 non-null    float64 
 1   sepal width (cm)   150 non-null    float64 
 2   petal length (cm)  150 non-null    float64 
 3   petal width (cm)   150 non-null    float64 
 4   species            150 non-null    category
dtypes: category(1), float64(4)
memory usage: 5.0 KB
None



Estadisticas descriptivas del conjunto de datos:
       sepal length (cm)  sepal width (cm)  petal length (cm)  petal width (cm)
count         150.000000        150.000000         150.000000        150.000000
mean            5.843333          3.057333           3.758000          1.199333
std             0.828066          0.435866           1.765298          0.762238
min             4.300000          2.000000           1.000000          0.100000
25%             5.100000          2.800000           1.600000          0.300000
50%             5.800000          3.000000           4.350000          1.300000
75%             6.400000          3.300000           5.100000          1.800000
max             7.900000          4.400000           6.900000          2.500000



Número de muestras por especie:
species
setosa        50
versicolor    50
virginica     50
Name: count, dtype: int64



Promedios de cada característica por especie:
            sepal length (cm)  sepal width (cm)  petal length (cm)  petal width (cm)
species                                                                             
setosa                  5.006             3.428              1.462             0.246
versicolor              5.936             2.770              4.260             1.326
virginica               6.588             2.974              5.552             2.026



Resumen estadístico agrupado por especie:
           sepal length (cm)                                              ... petal width (cm)                                          
                       count   mean       std  min    25%  50%  75%  max  ...            count   mean       std  min  25%  50%  75%  max
species                                                                   ...                                                           
setosa                  50.0  5.006  0.352490  4.3  4.800  5.0  5.2  5.8  ...             50.0  0.246  0.105386  0.1  0.2  0.2  0.3  0.6
versicolor              50.0  5.936  0.516171  4.9  5.600  5.9  6.3  7.0  ...             50.0  1.326  0.197753  1.0  1.2  1.3  1.5  1.8
virginica               50.0  6.588  0.635880  4.9  6.225  6.5  6.9  7.9  ...             50.0  2.026  0.274650  1.4  1.8  2.0  2.3  2.5

[3 rows x 32 columns]
<class 'pandas.DataFrame'>
RangeIndex: 150 entries, 0 to 149
Data columns (total 5 columns):
 #   Column             Non-Null Count  Dtype   
---  ------             --------------  -----   
 0   sepal length (cm)  150 non-null    float64 
 1   sepal width (cm)   150 non-null    float64 
 2   petal length (cm)  150 non-null    float64 
 3   petal width (cm)   150 non-null    float64 
 4   species            150 non-null    category
dtypes: category(1), float64(4)
memory usage: 5.0 KB
\end{consola}

\referencias

\end{document}